% The helper below gives readable standalone fallbacks for labels defined in
% Section 1.3. In the complete chapter, the real cleveref target is used.
\providecommand{\mextref}[2]{%
  \ifcsname r@#1\endcsname\cref{#1}\else#2\fi}

\section{Markov chains}
\label{m:sec:chains}

The preceding overview separated the stochastic processes used in this thesis
from the methods used to analyse them. Their common probabilistic structure is
the Markov property: conditional on the present state, the future is independent
of the past. This section introduces the relevant classes of process and records
the identities that make their behaviour concrete. Standard references are
Norris \cite{Norris1997}, Karlin and Taylor \cite{KarlinTaylor1975}, and, for
biological applications, Allen \cite{Allen2010}.

\subsection{Discrete-time Markov chains}
\label{m:sec:dtmc}

Let $\{X_n:n=0,1,2,\ldots\}$ take values in a countable state space
$\mathcal S$. It is a \emph{discrete-time Markov chain} if
\begin{equation}
\label{m:eq:dtmarkov}
\pr{X_{n+1}=j\mid X_n=i,X_{n-1},\ldots,X_0}
=\pr{X_{n+1}=j\mid X_n=i}=:P_{ij},
\end{equation}
where the right-hand side is independent of $n$ for a \emph{time-homogeneous}
chain. The transition matrix $P=(P_{ij})$ has non-negative entries and row sums
equal to one. Decomposing a path at an intermediate state gives the
Chapman--Kolmogorov identity
\begin{equation}
\label{m:eq:chapmankolmogorov}
\pr{X_{n+m}=j\mid X_0=i}
=\sum_{k\in\mathcal S}
  \pr{X_n=k\mid X_0=i}\,
  \pr{X_m=j\mid X_0=k},
\end{equation}
and hence $\pr{X_n=j\mid X_0=i}=(P^n)_{ij}$. Thus repeated steps correspond
to matrix multiplication. A state $i$ with $P_{ii}=1$ is \emph{absorbing}:
once entered, it cannot be left.

\subsubsection{Simple random walks}
\label{m:sec:randomwalk}

The simplest non-trivial chain moves one step at a time on the integers. Fix
$q\in(0,1)$ and set
\begin{equation}
\label{m:eq:rwrates}
P_{i,i+1}=q,
\qquad
P_{i,i-1}=1-q.
\end{equation}
Then $X_n=X_0+\sum_{k=1}^n\xi_k$, where the independent increments $\xi_k$
equal $+1$ with probability $q$ and $-1$ otherwise. The walk is symmetric
when $q=\tfrac12$, and, for deterministic $X_0$,
\begin{equation}
\label{m:eq:rwmoments}
\ex{X_n}=X_0+n(2q-1),
\qquad
\var{X_n}=4nq(1-q).
\end{equation}
The first quantity records the drift; the second records the diffusive spread
that the mean omits.

The full $n$-step law is also elementary. If $R$ of the $n$ increments are
positive, then $R\sim\operatorname{Bin}(n,q)$ and $X_n=X_0+2R-n$, so
\begin{equation}
\label{m:eq:rwlaw}
\pr{X_n=X_0+2r-n}
=\binom nr q^r(1-q)^{n-r},
\qquad r=0,1,\ldots,n.
\end{equation}
Only lattice points of the same parity as $X_0+n$ are possible
(\cref{m:fig:rwlaw}). This closed form depends on the increments being
independent of the current state. When the transition law depends on the
population already present, as it does for branching processes, generating
functions replace convolution; that method is developed in
\mextref{m:sec:pgfmethod}{Section~1.3.1.2}.

\begin{figure}[!htbp]
\centering
\includegraphics[width=0.70\textwidth]{rw_transition.pdf}
\caption{The displacement law \cref{m:eq:rwlaw} after $n=40$ steps, for
$q=0.4,0.5,0.6$ and $X_0=0$. Dotted lines mark the means $n(2q-1)$.
Changing $q$ mainly translates the distribution over this range, since its
standard deviation $\sqrt{4nq(1-q)}$ changes comparatively little.}
\label{m:fig:rwlaw}
\end{figure}

The random walk also appears inside continuous-time chains. Every such chain
separates into holding times and jumps, and the sequence of states visited at
the jumps is a discrete-time \emph{embedded jump chain}. For the linear
birth--death process of \cref{m:sec:bd}, the rates out of state $i\ge1$ are
$\lambda i$ and $\mu i$. Their common factor cancels, leaving the random walk
\cref{m:eq:rwrates} with
\begin{equation}
\label{m:eq:embedded}
q=\frac{\lambda}{\lambda+\mu}.
\end{equation}
Hitting probabilities therefore depend only on this embedded walk; holding
times matter when the question asks \emph{when} a state is reached.

For example, impose absorbing barriers at $0$ and $M$, and let $h_i$ be the
probability of reaching $M$ before $0$ from $i$. The first-step relation
$h_i=qh_{i+1}+(1-q)h_{i-1}$, with $h_0=0$ and $h_M=1$, gives
\begin{equation}
\label{m:eq:gamblersruin}
h_i=
\begin{cases}
\displaystyle\frac{1-\theta^i}{1-\theta^M},
& \displaystyle\theta:=\frac{1-q}{q}\ne1,\\[10pt]
\displaystyle\frac{i}{M},&q=\tfrac12.
\end{cases}
\end{equation}
This gambler's-ruin calculation is the elementary instance of the general
first-step principle in \mextref{m:sec:firststep}{Section~1.3.1.1}.

\begin{figure}[!htbp]
\centering
\includegraphics[width=0.94\textwidth]{random_walk.pdf}
\caption{Two views of the walk \cref{m:eq:rwrates}. (a) Seven paths from
$X_0=0$ with $q=0.45$, together with the mean and one-standard-deviation
envelope from \cref{m:eq:rwmoments}. (b) The probability
\cref{m:eq:gamblersruin} of reaching $M=20$ before $0$. The symmetric case
$\theta=1$ is linear; drift towards $M$ ($\theta<1$) bends the probability
upwards, while drift towards $0$ bends it downwards.}
\label{m:fig:randomwalk}
\end{figure}

\subsubsection{Galton--Watson branching processes}
\label{m:sec:gw}

The Galton--Watson process is a chain on the non-negative integers in which
each member of one generation independently produces the next. It is the
classical model of stochastic branching and the principal discrete-time
process in this thesis; standard accounts are Harris \cite{Harris1963} and
Athreya and Ney \cite{Athreya1972}.

Let $L$ be the offspring count of one individual, with
$\pr{L=k}=p_k$ and mean
\begin{equation*}
m:=\ex{L}.
\end{equation*}
Starting from one progenitor, $Z_0=1$, the population evolves as
\begin{equation}
\label{m:eq:gwrecursion}
Z_{n+1}=\sum_{i=1}^{Z_n}L_i,
\end{equation}
where the $L_i$ are independent copies of $L$. Conditioning on $Z_n$ gives
$\ex{Z_{n+1}}=m\ex{Z_n}$ and therefore
$\ex{Z_n}=m^n$. The regimes $m<1$, $m=1$, and $m>1$ are called
subcritical, critical, and supercritical.

The random sum in \cref{m:eq:gwrecursion} becomes a composition under a
\pgf{}. Write $\phi(z)=\ex{z^L}$ and let $\phi_n$ be the generating function
of $Z_n$. Then
\begin{equation}
\label{m:eq:gwcomposition}
\phi_{n+1}=\phi_n\circ\phi,
\qquad \phi_0(z)=z,
\end{equation}
so $\phi_n$ is the $n$-fold iterate of the offspring generating function.
The process is therefore encoded by the orbit of a single map, although that
orbit seldom has a closed form.

\paragraph{The binary case.}
The thesis specialises to the binary, or \emph{simple}, process in which an
individual divides into two or dies:
\begin{equation}
\label{m:eq:binarylaw}
p_2=p\quad\text{(division)},
\qquad
p_0=1-p\quad\text{(death)},
\qquad
p_k=0\quad(k\notin\{0,2\}).
\end{equation}
Here $m=2p$, $\ex{Z_n}=(2p)^n$, and
\begin{equation}
\label{m:eq:pgfbinary}
\phi(z)=pz^2+(1-p).
\end{equation}
If $u_n=\pr{Z_n=0}$ and $S_n=1-u_n=\pr{Z_n>0}$, evaluating
\cref{m:eq:gwcomposition} at $z=0$ gives
\begin{equation}
\label{m:eq:SnIter}
S_{n+1}=2pS_n-pS_n^2,
\qquad S_0=1.
\end{equation}
For $0<p\le\tfrac12$, $S_n$ decreases to zero but is positive at every finite
generation; for $p>\tfrac12$ it decreases to the positive fixed point
$(2p-1)/p$. Thus subcritical and critical survival generally \emph{converges}
to zero rather than reaching it after finitely many iterations. The following
chapter studies the rate of this convergence and the constant that remains
after conditioning on survival; the present chapter needs only the process,
its mean, and the functional iteration above.

\FloatBarrier

\subsection{Continuous-time Markov chains}
\label{m:sec:ctmc}

Let $\{X_t:t\ge0\}$ take values in a countable state space $\mathcal S$. It
is a \emph{continuous-time Markov chain} if, for all $s,t\ge0$,
\begin{equation}
\label{m:eq:markov}
\pr{X_{t+s}=j\mid X_t=i,\{X_u\}_{u<t}}
=\pr{X_{t+s}=j\mid X_t=i},
\end{equation}
and it is time homogeneous when the right-hand side depends on $s$ but not
on $t$. A homogeneous chain is specified by its generator $Q=(q_{ij})$,
whose off-diagonal entries are transition rates:
\begin{equation}
\label{m:eq:generator}
q_{ij}=\lim_{\Delta t\downarrow0}
\frac{\pr{X_{t+\Delta t}=j\mid X_t=i}}{\Delta t}\quad(j\ne i),
\qquad
q_{ii}=-\sum_{j\ne i}q_{ij}=:-q_i.
\end{equation}
Equivalently,
\begin{equation}
\label{m:eq:infinitesimal}
\pr{X_{t+\Delta t}=j\mid X_t=i}
=q_{ij}\Delta t+o(\Delta t)\qquad(j\ne i),
\end{equation}
with probability $1-q_i\Delta t+o(\Delta t)$ of no jump. A state with
$q_i=0$ is absorbing.

\subsubsection{Exponential clocks}
\label{m:sec:exp}

The holding times of a homogeneous chain are exponential. A non-negative
\rv{} $T$ has the exponential distribution with rate $\theta>0$, written
$T\sim\operatorname{Exp}(\theta)$, when
\begin{equation}
\label{m:eq:expdef}
f_T(t)=\theta\ee^{-\theta t},
\qquad
\pr{T>t}=\ee^{-\theta t}
\qquad(t\ge0).
\end{equation}
Its mean is $1/\theta$ and its variance is $1/\theta^2$. The decisive
property is memorylessness:
\begin{equation}
\label{m:eq:memoryless}
\pr{T>t_1+t_2\mid T>t_1}
=\frac{\ee^{-\theta(t_1+t_2)}}{\ee^{-\theta t_1}}
=\pr{T>t_2}.
\end{equation}
Thus the elapsed holding time carries no information about the residual
wait. Among non-degenerate continuous waiting-time distributions on the
non-negative line, this property characterises the exponential law.

Several possible transitions can be represented by independent clocks.

\begin{proposition}[Competing clocks]
\label{m:prop:race}
Let $T_1,\ldots,T_r$ be independent, with
$T_k\sim\operatorname{Exp}(\theta_k)$. If
$T=\min_kT_k$ and $\Theta=\sum_k\theta_k$, then
\begin{equation}
\label{m:eq:race}
T\sim\operatorname{Exp}(\Theta),
\qquad
\pr{T_k=T}=\frac{\theta_k}{\Theta},
\end{equation}
and the winning clock is independent of $T$.
\end{proposition}

\begin{proof}
Independence gives
$\pr{T>t}=\prod_k\ee^{-\theta_kt}=\ee^{-\Theta t}$. Moreover,
\[
\pr{T_k=T,\ T>t}
=\int_t^\infty\theta_k\ee^{-\Theta u}\,\dd u
=\frac{\theta_k}{\Theta}\ee^{-\Theta t},
\]
which gives both the probability of clock $k$ winning and its independence
from the minimum.
\end{proof}

Consequently, a chain leaves state $i$ after an
$\operatorname{Exp}(q_i)$ holding time and then jumps to $j$ with
probability $q_{ij}/q_i$, independently of the holding time. This separates
the timing of a jump from its destination and produces the embedded chain
used in \cref{m:eq:embedded}.

\begin{remark}[Direct simulation]
\label{m:rem:gillespie}
The same construction gives Gillespie's direct method \cite{Gillespie1977}.
At the current state, calculate the event rates $\theta_k$ and their total
$\Theta$; draw a holding time from $\operatorname{Exp}(\Theta)$; select event
$k$ with probability $\theta_k/\Theta$; update the state and repeat. The
method generates exact CTMC paths, up to numerical round-off, and is used in
this thesis to check analytical calculations and to study models without
closed transient laws.
\end{remark}

\subsubsection{The Poisson process}
\label{m:sec:poisson}

The Poisson process counts events that occur at a constant rate. Let $N_t$
take values in $\{0,1,2,\ldots\}$, with the sole non-zero transition rate
\begin{equation}
\label{m:eq:poissonrates}
q_{n,n+1}=\alpha,
\qquad n\ge0.
\end{equation}
Every holding time is therefore an independent
$\operatorname{Exp}(\alpha)$ variable. Starting from $N_0=0$, the count at
time $t$ has law
\begin{equation}
\label{m:eq:poissonlaw}
\pr{N_t=n}=\ee^{-\alpha t}\frac{(\alpha t)^n}{n!},
\qquad
\ex{N_t}=\var{N_t}=\alpha t.
\end{equation}
The equality of mean and variance is a characteristic baseline for comparison
with the over-dispersion of branching populations.

\begin{figure}[!htbp]
\centering
\includegraphics[width=0.94\textwidth]{poisson_process.pdf}
\caption{A Poisson process of rate $\alpha=1.5$. (a) One path, with its
exponential holding intervals and mean $\alpha t$. (b) The law
\cref{m:eq:poissonlaw} at $t=1,3,6$. Its absolute width grows with time, while
its coefficient of variation $1/\sqrt{\alpha t}$ decreases.}
\label{m:fig:poisson}
\end{figure}

Counts in disjoint intervals are independent, and the distribution of an
increment depends only on the interval length. If each event is retained
independently with probability $\varpi$, the retained events form a Poisson
process of rate $\varpi\alpha$; this thinning property will also be useful for
simulating time-dependent rates.

\subsubsection{The Yule process}
\label{m:sec:yule}

Replacing the constant Poisson rate by a rate proportional to the population
produces the \emph{Yule process}. Each individual gives birth at rate
$\lambda$ and no individual dies, so, from a single founder, the only non-zero
rate is
\begin{equation}
\label{m:eq:yulerates}
q_{n,n+1}=\lambda n,
\qquad n\ge1.
\end{equation}
The holding time in state $n$ is $\operatorname{Exp}(\lambda n)$ and shortens
as the population grows. The count $Y_t$ is geometric:
\begin{equation}
\label{m:eq:yulelaw}
\pr{Y_t=n}=\ee^{-\lambda t}
\lrb{1-\ee^{-\lambda t}}^{n-1},
\qquad
\ex{Y_t}=\ee^{\lambda t},
\qquad
\var{Y_t}=\ee^{\lambda t}\lrb{\ee^{\lambda t}-1},
\end{equation}
for $n=1,2,\ldots$. Unlike the Poisson process, its relative fluctuations do
not vanish: the coefficient of variation tends to one. Early random
differences are amplified by later births. A process begun from $k$ founders
is the sum of $k$ independent one-founder processes and has the corresponding
negative-binomial law.

\begin{remark}[Origin and purpose]
\label{m:rem:yuleorigin}
The process arose from a biological question rather than as an abstract
chain. Willis argued in \emph{Age and Area} \cite{Willis1922} that the highly
skewed distribution of species among genera required a quantitative account.
Yule supplied one in 1925 \cite{Yule1925}: species arise within a genus by a
neutral pure-birth mechanism of the form above, while new genera continue to
appear. Mixing the geometric within-genus law over genus ages produces the
heavy-tailed genus-size frequencies later associated with the Yule--Simon
distribution, after Simon's urn-model treatment \cite{Simon1955}. The same
linear pure-birth dynamics appeared independently in Furry's 1937 model of
cosmic-ray cascades \cite{Furry1937}, hence the name Yule--Furry process in
the physics literature. The subsequent probability theory of birth--death
processes, including Kendall's 1948 treatment \cite{Kendall1948}, developed
largely through Feller's and Bartlett's line of work rather than from Yule's
paper. The model now serves as the null process of diversification without
extinction: its pure-birth limit underlies reconstructed phylogenetic
likelihoods \cite{Nee1994} and the one-parameter Yule tree prior implemented
in Bayesian software such as BEAST \cite{Bouckaert2014}.
\end{remark}

\subsubsection{Birth--death processes}
\label{m:sec:bd}

Let $X_t$ count individuals that independently give birth at rate $\lambda$
and die at rate $\mu$. The non-zero rates out of $n\ge1$ are
\begin{equation}
\label{m:eq:bdrates}
q_{n,n+1}=\lambda n,
\qquad
q_{n,n-1}=\mu n.
\end{equation}
State $0$ is absorbing, and $\mu=0$ recovers the Yule process. This
\emph{linear birth--death process} is the continuous-time counterpart of the
Galton--Watson process and the running example for the methods of
\mextref{m:sec:ctmethods}{Section~1.3.2}. From $X_0=N$, its mean is
$N\ee^{(\lambda-\mu)t}$. Its embedded chain is \cref{m:eq:rwrates} with
$q=\lambda/(\lambda+\mu)$, so the gambler's-ruin calculation determines
which boundaries are reached even though it does not determine when.

The comparison between paths and their mean is particularly instructive.

\begin{figure}[!htbp]
\centering
\includegraphics[width=0.94\textwidth]{birth_death_paths.pdf}
\caption{Nine paths of the linear birth--death process from $X_0=5$, with
$\mu=1$ and three values of $\lambda$; red curves show
$5\ee^{(\lambda-\mu)t}$. At criticality the mean remains five although some
paths are already extinct and others are much larger. Supercritical growth
also permits early extinction, whose eventual probability here is
$(\mu/\lambda)^5$.}
\label{m:fig:bdpaths}
\end{figure}

For one founder, extinction is certain when $\lambda\le\mu$ and has probability
$\mu/\lambda$ when $\lambda>\mu$; from $N$ founders the latter probability is
$(\mu/\lambda)^N$. Thus even a correct mean can conceal the absorbing fate of
most or some of the population. General birth--death chains permit
state-dependent rates $\lambda_n$ and $\mu_n$; the linear case is special
because individuals act independently and its generating functions and
moments close.

\FloatBarrier

\subsubsection{Birth--death--catastrophe processes}
\label{m:sec:bdc}

The process central to the later modelling chapters adds a third event. In
addition to individual births and deaths, any replicative unit may trigger an
abrupt catastrophe that terminates the whole intracellular population.

\begin{definition}[Birth--death--catastrophe process]
\label{m:def:bdc}
Let $X_t$ take values in $\{0,1,2,\ldots\}\cup\{H\}$ and record the
intracellular count while the compartment remains intact. Each replicative
unit independently gives birth at rate $\lambda$, dies at rate $\mu$, and
triggers catastrophe at rate $\delta$. For $n\ge1$, the non-zero transition
rates are
\begin{equation}
\label{m:eq:bdcrates}
q_{n,n+1}=\lambda n,
\qquad
q_{n,n-1}=\mu n,
\qquad
q_{n,H}=\delta n,
\end{equation}
and both $0$ and $H$ are absorbing.
\end{definition}

The absorbing states represent different terminal fates. State $0$ is
internal extinction after the last ordinary death; $H$ is catastrophe or
rupture and must not be identified with $0$. Let
$\tau=\inf\{t:X_t=H\}$ be the catastrophe time. The released-count process
used in Chapter~3 is
\begin{equation}
\label{m:eq:bdcrelease}
W_t=
\begin{cases}
0,&t<\tau,\\
X_{\tau-},&t\ge\tau,
\end{cases}
\qquad
K:=X_{\tau-},
\end{equation}
where $K$ is the burst size on paths that rupture. The left limit is essential:
at time $\tau$ the internal state jumps to $H$, whereas the released load is
the count present immediately beforehand.

Because the total catastrophe rate is $\delta n$ in state $n$, rupture is more
likely from a fuller compartment. Before the first trigger, the family trees
and catastrophe clocks descending from distinct founders can be constructed
independently. If
$I(t)=\pr{X_t\ne H\mid X_0=1}$ is the non-catastrophe probability, then the
corresponding probability from $k$ founders is $I(t)^k$. This factorisation
does not make post-trigger lineages independent: the first catastrophe stops
them all simultaneously.

\begin{figure}[!htbp]
\centering
\resizebox{0.92\textwidth}{!}{%
\begin{tikzpicture}[
  >={Stealth[length=1.9mm]},
  st/.style={circle,draw,thick,minimum size=7.5mm,inner sep=0pt,font=\small},
  ab/.style={circle,draw,thick,double,double distance=0.7pt,
             minimum size=7.5mm,inner sep=0pt,font=\small},
  lb/.style={font=\footnotesize,inner sep=1.6pt},
  rowlab/.style={font=\small\itshape,anchor=west},
  every path/.style={thick}
]
\node[rowlab] at (-1.9,1.15) {birth--death};
\node[ab] (a0) at (0,0) {$0$};
\node[st] (a1) at (2.1,0) {$1$};
\node[st] (a2) at (4.2,0) {$2$};
\node[st] (a3) at (6.3,0) {$3$};
\node (a4) at (7.9,0) {$\cdots$};
\draw[->] (a1) to[bend left=38] node[above,lb]{$\lambda$} (a2);
\draw[->] (a2) to[bend left=38] node[above,lb]{$2\lambda$} (a3);
\draw[->] (a3) to[bend left=38] node[above,lb]{$3\lambda$} (a4);
\draw[->] (a1) to[bend left=38] node[below,lb]{$\mu$} (a0);
\draw[->] (a2) to[bend left=38] node[below,lb]{$2\mu$} (a1);
\draw[->] (a3) to[bend left=38] node[below,lb]{$3\mu$} (a2);

\begin{scope}[yshift=-3.5cm]
\node[rowlab] at (-1.9,1.15) {birth--death--catastrophe};
\node[ab] (b0) at (0,0) {$0$};
\node[st] (b1) at (2.1,0) {$1$};
\node[st] (b2) at (4.2,0) {$2$};
\node[st] (b3) at (6.3,0) {$3$};
\node (b4) at (7.9,0) {$\cdots$};
\draw[->] (b1) to[bend left=38] node[above,lb]{$\lambda$} (b2);
\draw[->] (b2) to[bend left=38] node[above,lb]{$2\lambda$} (b3);
\draw[->] (b3) to[bend left=38] node[above,lb]{$3\lambda$} (b4);
\draw[->] (b1) to[bend left=38] node[below,lb]{$\mu$} (b0);
\draw[->] (b2) to[bend left=38] node[below,lb]{$2\mu$} (b1);
\draw[->] (b3) to[bend left=38] node[below,lb]{$3\mu$} (b2);
\node[ab] (bh) at (4.2,-2.35) {$H$};
\draw[->,densely dashed] (b1) to[bend right=12]
  node[pos=0.68,left=1.5pt,lb]{$\delta$} (bh);
\draw[->,densely dashed] (b2) --
  node[pos=0.68,right=1.5pt,lb]{$2\delta$} (bh);
\draw[->,densely dashed] (b3) to[bend left=12]
  node[pos=0.68,right=1.5pt,lb]{$3\delta$} (bh);
\end{scope}
\end{tikzpicture}}
\caption{Linear birth--death dynamics (top) and the
birth--death--catastrophe extension (bottom). Double circles are absorbing.
Ordinary extinction enters $0$ only from state $1$; catastrophe instead jumps
from any positive count to the distinct state $H$, at total rate $\delta n$.}
\label{m:fig:ratediagram}
\end{figure}

The detailed BDC analysis is deferred to
\fwd{bdc}{the chapter on the birth--death--catastrophe process}, where
non-catastrophe, internal-extinction and non-fixation probabilities are kept
separate and the moments of $X_t$ and $W_t$ are derived. One modelling caution
is important enough to establish here before making that deferral.

\begin{remark}[A state-dependent catastrophe is not a mean-field hazard]
\label{m:rem:notmeanfield}
Write $\chi\in(0,1)$ for the probability that one individual initiates
catastrophe during a discrete generation. This is a probability per
generation, distinct from the continuous-time per-capita rate $\delta$. Let
$\widetilde X_i$ be the branching population with catastrophe disabled and
let its cumulative exposure before generation $n$ be
\[
\mathcal E_n=\sum_{i=0}^{n-1}\widetilde X_i.
\]
Conditional on this catastrophe-free genealogy, every exposed individual
fails to trigger independently, and therefore
\begin{equation}
\label{m:eq:norupture}
\pr{\text{no catastrophe before generation }n}
=\ex{(1-\chi)^{\mathcal E_n}}.
\end{equation}
Replacing $\mathcal E_n$ by its mean is not exact. Since
$e\mapsto(1-\chi)^e$ is convex, Jensen's inequality gives
\[
\ex{(1-\chi)^{\mathcal E_n}}
\ge(1-\chi)^{\ex{\mathcal E_n}},
\]
so the mean-population substitution underestimates the non-catastrophe
probability and overestimates catastrophe risk. Nor is non-catastrophe the
same as continued population survival: their intersection has probability
$\ex{\ind\{\widetilde X_n>0\}(1-\chi)^{\mathcal E_n}}$, to which this Jensen
bound does not apply. The exact hazard depends on cumulative path exposure,
not on the population mean at a single time.
\end{remark}

\FloatBarrier

\subsection{Time-inhomogeneous processes}
\label{m:sec:inhomogeneous}

The preceding chains are time homogeneous: their rates depend on the state
but not on the clock. A time-inhomogeneous CTMC has a generator $Q(t)$ and
two-time transition probabilities
\begin{equation}
\label{m:eq:inhomogtrans}
p_{ij}(s,t)=\pr{X_t=j\mid X_s=i},
\qquad s\le t.
\end{equation}
For $s\le u\le t$, Chapman--Kolmogorov becomes
$P(s,t)=P(s,u)P(u,t)$. Under the usual regularity conditions, differentiation
at either endpoint gives
\begin{equation}
\label{m:eq:inhomogkolmogorov}
\pdd{}{t}P(s,t)=P(s,t)Q(t),
\qquad
\pdd{}{s}P(s,t)=-Q(s)P(s,t).
\end{equation}
The solution is generally a time-ordered product rather than
$\ee^{(t-s)Q}$; the ordinary exponential is recovered when the generators
at different times commute. Simulation likewise uses integrated hazards or
thinning, since event rates can change between jumps even while the state is
fixed.

\paragraph{Linear birth--death rates.}
\label{m:sec:tibd}
Suppose the per-capita birth and death rates are prescribed calendar-time
functions $\lambda(t)$ and $\mu(t)$. Linearity still closes the mean:
\begin{equation}
\label{m:eq:meanodeinhom}
\ex{X_t\mid X_0=n_0}
=n_0\exp\lrb{\int_0^t\bigl(\lambda(r)-\mu(r)\bigr)\,\dd r}.
\end{equation}
Extinction requires more care because the starting and terminal times play
different roles. The natural backward quantity is
\[
u(s,T):=\pr{X_T=0\mid X_s=1}.
\]
A first-step expansion at time $s$ gives the backward Riccati equation
\begin{equation}
\label{m:eq:riccatiext}
-\pdd{u}{s}
=\mu(s)-\bigl(\lambda(s)+\mu(s)\bigr)u+\lambda(s)u^2,
\qquad
u(T,T)=0.
\end{equation}
By contrast, the forward-endpoint probability
$u_{\mathrm f}(t):=\pr{X_t=0\mid X_0=1}=u(0,t)$ is, in Kendall's notation
\cite{Kendall1948},
\begin{equation}
\label{m:eq:inhomextforward}
u_{\mathrm f}(t)=\frac{B(t)}{1+B(t)},
\qquad
B(t)=\int_0^t\mu(s)
\exp\lrb{\int_0^s\bigl(\mu(r)-\lambda(r)\bigr)\,\dd r}\,\dd s.
\end{equation}
It does not generally satisfy \cref{m:eq:riccatiext} with $s$ replaced by
$t$. If time is reparametrised by the remaining horizon,
$v(r)=u(T-r,T)$, then the Riccati coefficients are
$\lambda(T-r)$ and $\mu(T-r)$; the familiar one-variable homogeneous formula
is recovered when the rates are constant. This distinction prevents a
calendar-time rate schedule from being read backwards inadvertently.

\subsubsection{The seasonal Lotka--Volterra process}
\label{m:sec:seasonallv}

Externally imposed time dependence remains visible when populations interact.
Let $X_t$ and $Y_t$ be prey and predator counts. Prey reproduction is seasonal,
predation converts one prey into one predator, and predators die individually:
\begin{equation}
\label{m:eq:lvrate}
q_{(x,y),(x+1,y)}(t)=b\mathcal S^+(t)x,
\qquad
q_{(x,y),(x-1,y+1)}(t)=exy,
\qquad
q_{(x,y),(x,y-1)}(t)=dy.
\end{equation}
The seasonal multiplier is the positive part of a sinusoid,
\begin{equation}
\label{m:eq:lvkernel}
\mathcal S^+(t)=\max\{0,1+A\sin(\omega t)\}.
\end{equation}
Clipping is required because a transition rate cannot be negative. For $A>1$,
each period $2\pi/\omega$ contains a closed season of length
$(2/\omega)\arccos(1/A)$ during which prey births stop.

On the predator-free axis, $Y_t\equiv0$ and $X_t$ is a time-inhomogeneous
Yule process with per-capita rate $b\mathcal S^+(t)$. Hence
\[
\ex{X_t\mid X_0=x_0,Y_0=0}
=x_0\exp\lrb{b\int_0^t\mathcal S^+(s)\,\dd s},
\]
and extinction is impossible on that axis. In the interior, the interaction
is quadratic. The exact first moments involve $\ex{X_tY_t}$ and do not close;
for comparison, the corresponding deterministic system is
\begin{equation}
\label{m:eq:lvmeanfield}
\dot x=b\mathcal S^+(t)x-exy,
\qquad
\dot y=exy-dy.
\end{equation}
Thus this model marks the limit of what a closed mean can reveal: the forcing
is prescribed, but the interacting stochastic law is not reduced to
\cref{m:eq:meanodeinhom}. Exact paths can be simulated by thinning against
$\mathcal S^+(t)\le1+A$.

\begin{figure}[!htbp]
\centering
\includegraphics[width=0.78\textwidth]{seasonal_lv.pdf}
\caption{The seasonal Lotka--Volterra process with $A=2$, $\omega=2\pi$,
$b=2$, $e=0.005$, $d=2$, and $(X_0,Y_0)=(300,100)$. (a) Clipping the forcing
creates a closed season occupying one third of every period. (b) Ten exact
prey--predator paths. (c) Pointwise medians and 10--90\% bands from 200 paths,
with \cref{m:eq:lvmeanfield} dashed. The first cycle is calendar-locked and
tightly concentrated; later interaction amplifies phase and amplitude
differences. If prey reach zero, the remaining predators form a pure-death
process and the path eventually reaches $(0,0)$.}
\label{m:fig:seasonallv}
\end{figure}

\FloatBarrier

\subsubsection{The logistic speciation model}
\label{m:sec:logisticspeciation}

The second example is a pure-birth species process driven by a prescribed
deterministic environment. Let $N(t)$ be total abundance across all species,
assumed large enough to follow the logistic equation
\begin{equation}
\label{m:eq:logisticode}
\dda{N}{t}=rN\lrb{1-\frac{N}{C}},
\qquad
N(0)=N_0,
\qquad 0<N_0<C,
\end{equation}
where $C$ is environmental capacity. Its solution is
\begin{equation}
\label{m:eq:logisticsol}
N(t)=\frac{C\ee^{rt}}{\xi+\ee^{rt}},
\qquad
\xi=\frac{C}{N_0}-1.
\end{equation}

Let $S_t$ be the number of species, with $S_0=1$. Species do not become extinct
in this basic model. Each existing species gives rise to a new one at the
per-species rate
\begin{equation}
\label{m:eq:speciationrate}
\beta(t)=\sigma\bigl(C-N(t)\bigr),
\end{equation}
so $q_{s,s+1}(t)=\beta(t)s$. The model is therefore exactly a Yule process
under a deterministic time change; it is not claimed to be an exact reduction
of a state-dependent stochastic population model. Substitution of
\cref{m:eq:logisticsol} gives
\begin{equation}
\label{m:eq:lambdat}
\beta(t)=\sigma C\frac{\xi}{\xi+\ee^{rt}},
\end{equation}
which decreases to zero as the environment fills. Its cumulative intensity is
\begin{equation}
\label{m:eq:lambdaint}
\mathcal B(t):=\int_0^t\beta(s)\,\dd s
=\sigma Ct+\frac{\sigma C}{r}
\log\lrb{\frac{\xi+1}{\xi+\ee^{rt}}}.
\end{equation}
Replacing $\lambda t$ by $\mathcal B(t)$ in the Yule law
\cref{m:eq:yulelaw} gives the complete fixed-time distribution. In particular,
\begin{equation}
\label{m:eq:logspecmean}
\ex{S_t}=\ee^{\mathcal B(t)}
=\ee^{\sigma Ct}
\lrb{\frac{\xi+1}{\xi+\ee^{rt}}}^{\sigma C/r}
\xrightarrow[t\to\infty]{}
(\xi+1)^{\sigma C/r}
=\lrb{\frac{C}{N_0}}^{\sigma C/r}.
\end{equation}
The count is non-decreasing and has a bounded limiting mean, so it converges
almost surely to a finite random species count. This is a one-way coupling:
the environment sets the speciation clock, while $S_t$ does not feed back into
$N(t)$.

\begin{figure}[!htbp]
\centering
\includegraphics[width=0.73\textwidth]{logspec_yule_mean.pdf}
\caption{Mean species count \cref{m:eq:logspecmean} for
$N_0=1,2,5$, with $C=10$, $r=1$, and $\sigma=0.08$; dashed lines are the
limits $(C/N_0)^{\sigma C/r}$. A smaller founding abundance leaves more
integrated unused capacity, so the limiting mean is higher.}
\label{m:fig:logspecmean}
\end{figure}

The same deterministic logistic environment can drive nearby stochastic
models without changing the status of the process just analysed. In
\cref{m:fig:lsmean}, the per-species origination rate is instead
$\beta_{\mathrm{bd}}(t)=\sigma N'(t)$ and each species becomes extinct at a
constant per-species rate $\varepsilon>0$. Its mean can therefore rise and
then fall. This is a numerical comparison, not an alternative interpretation
of \cref{m:eq:speciationrate}: the basic model above has no species extinction
and its paths remain non-decreasing.

\begin{figure}[!htbp]
\centering
\includegraphics[width=0.58\textwidth]{LSMean.png}
\caption{Numerical comparison with a logistic-environment birth--death
variant. Pale blue curves are species-count paths and the dark blue curve is
their sample mean; the lower panel shows the per-species origination rate
$\sigma N'(t)$ and a constant extinction rate. Richness turns over after
origination falls below extinction. The ordinate $X_t$ in the supplied
illustration denotes the species count $S_t$ used in this section; this is not
the pure-birth model of \cref{m:eq:speciationrate,m:eq:logspecmean}.}
\label{m:fig:lsmean}
\end{figure}

\FloatBarrier

\subsection{Coupled ODE--CTMC systems}
\label{m:sec:coupled}

A continuous trajectory and a discrete chain can be combined in three
distinct ways. Write $\mathbf y(t)\in\mathbb R^d$ for the continuous state and
$M_t\in\mathcal S$ for the discrete mode.

\begin{enumerate}[label=(\roman*),leftmargin=*,itemsep=0.45em]
\item An autonomous Markov chain may switch the vector field:
\begin{equation}
\label{m:eq:coupledone}
\dda{\mathbf y}{t}=F\bigl(\mathbf y(t),M_t\bigr),
\qquad M_t\text{ autonomous}.
\end{equation}
Here $M_t$ remains Markov on its own; $\mathbf y(t)$ generally does not,
because its future also depends on the unobserved current mode.

\item A prescribed continuous trajectory may drive the chain rates:
\begin{equation}
\label{m:eq:coupledtwo}
\dda{\mathbf y}{t}=F\bigl(\mathbf y(t)\bigr),
\qquad
q_{ij}(t)=q_{ij}\bigl(\mathbf y(t)\bigr).
\end{equation}
The trajectory can be solved first, and $M_t$ is then a time-inhomogeneous
Markov chain. The logistic speciation model is of this form.

\item In genuine two-way coupling, each component drives the other:
\begin{equation}
\label{m:eq:coupledthree}
\dda{\mathbf y}{t}=F\bigl(\mathbf y(t),M_t\bigr),
\qquad
q_{ij}=q_{ij}\bigl(\mathbf y(t)\bigr).
\end{equation}
Neither component can generally be solved in isolation.
\end{enumerate}

In all three cases the pair is naturally a piecewise-deterministic Markov
process \cite{Davis1993}; in the two-way case it is generally the pair, rather
than either component alone, that is Markov. For a suitable test function
$f(\mathbf y,i)$, the joint generator is
\begin{equation}
\label{m:eq:coupledgenerator}
(\mathcal Lf)(\mathbf y,i)
=F(\mathbf y,i)\mathbin{\cdot}\nabla_{\mathbf y}f(\mathbf y,i)
+\sum_{j\ne i}q_{ij}(\mathbf y)
  \bigl[f(\mathbf y,j)-f(\mathbf y,i)\bigr].
\end{equation}
The first term transports the continuous state along its deterministic flow;
the second is the usual jump generator. Removing the $\mathbf y$-dependence
of the rates gives case (i), while removing the mode-dependence of $F$ gives
case (ii).

\subsubsection{Accumulation with gated release}
\label{m:sec:gated}

Let $A_t$ be an accumulated amount, $R_t$ the amount transferred into a
receiving population, and $G_t\in\{0,1\}$ a closed/open gate. Between gate
jumps,
\begin{equation}
\label{m:eq:gatedflow}
\dda{A_t}{t}=\rho-G_tcA_t,
\qquad
\dda{R_t}{t}=G_tcA_t,
\end{equation}
where $\rho>0$ is the external input and $c>0$ the transfer rate. The gate
opens faster when more has accumulated and closes at a constant rate:
\begin{equation}
\label{m:eq:gatedjumps}
\pr{G_{t+\Delta t}=1\mid G_t=0,A_t=a}
=\eta a\,\Delta t+o(\Delta t),
\qquad
\pr{G_{t+\Delta t}=0\mid G_t=1}
=\gamma\,\Delta t+o(\Delta t).
\end{equation}
Thus $\eta$ has units of inverse amount--time, whereas $\gamma$ is a rate.
Both directions of \cref{m:eq:coupledthree} are present.

Adding the flow equations gives the pathwise balance law
\begin{equation}
\label{m:eq:gatedbalance}
A_t+R_t=A_0+R_0+\rho t.
\end{equation}
Transfer conserves the total, and only the external input changes it. From
non-negative initial conditions, both components remain non-negative and
$R_t$ is non-decreasing. The first moment already exposes the coupling:
$\dda{}{t}\ex{A_t}=\rho-c\ex{G_tA_t}$, which is not closed in
$\ex{A_t}$ and $\ex{G_t}$ separately.

\begin{figure}[!htbp]
\centering
\includegraphics[width=0.68\textwidth]{gated_paths.pdf}
\caption{Three exact paths of gated release from
$(A_0,R_0,G_0)=(5,0,0)$, with $\rho=2$, $c=4$, $\eta=0.015$, and
$\gamma=1.45$. While the gate is closed, $A_t$ grows linearly and $R_t$ is
constant; an open gate transfers the accumulated amount into $R_t$ and pulls
$A_t$ towards $\rho/c$. State-dependent opening produces intermittent
releases.}
\label{m:fig:gatedpaths}
\end{figure}

\FloatBarrier

\subsubsection{Accumulation with surging gated release}
\label{m:sec:surging}

The surging variant retains the flow \cref{m:eq:gatedflow} and closing rate
$\gamma$, but opens the gate autonomously at rate $\nu$:
\begin{equation}
\label{m:eq:surgingjumps}
\pr{G_{t+\Delta t}=1\mid G_t=0}
=\nu\,\Delta t+o(\Delta t).
\end{equation}
Here $G_t$ is an autonomous two-state CTMC, so the model is case (i) of the
taxonomy. Its mean closes:
\begin{equation}
\label{m:eq:gatemean}
\dda{}{t}\ex{G_t}
=\nu\bigl(1-\ex{G_t}\bigr)-\gamma\ex{G_t},
\qquad
\ex{G_t}\longrightarrow\frac{\nu}{\nu+\gamma}.
\end{equation}
The population moment still contains $\ex{G_tA_t}$, since the flow is driven
by the gate. The balance law \cref{m:eq:gatedbalance} is unchanged.

\begin{figure}[!htbp]
\centering
\includegraphics[width=0.94\textwidth]{gated_ensemble.pdf}
\caption{Medians and pointwise 10--90\% bands from 2000 paths. Both models
use $\rho=2$, $c=4$, $\gamma=1.45$, and $(A_0,R_0,G_0)=(5,0,0)$. Gated
release (left) opens at intensity $\eta A_t$ with $\eta=0.015$; surging
release (right) opens at the autonomous rate $\nu=0.3$. The parameters have
different dimensions, and the two opening mechanisms produce visibly
different accumulation levels.}
\label{m:fig:gatedensemble}
\end{figure}

\subsubsection{Simulating a coupled system}
\label{m:sec:coupledsim}

The competing-clock construction extends to the joint process, but the
continuous motion makes the jump rates vary between events. Suppose the
current absolute time is $t$, the discrete mode is $i$, and
$\mathbf y_u$ denotes the deterministic flow with that mode held fixed. If
$T$ is the next jump time, then, for $s\ge t$,
\begin{equation}
\label{m:eq:coupledholding}
\pr{T>s\mid\mathbf y_t,M_t=i}
=\exp\lrb{-\int_t^s q_i(\mathbf y_u)\,\dd u},
\qquad
q_i(\mathbf y)=\sum_{j\ne i}q_{ij}(\mathbf y).
\end{equation}
At $T$, the destination is selected with probability
$q_{ij}(\mathbf y_T)/q_i(\mathbf y_T)$. The integrated hazard can be inverted
directly when a closed form is available, or handled by quadrature or
thinning otherwise.

For gated release, a closed gate gives $A_{t+r}=A_t+\rho r$, so its cumulative
opening hazard over a wait $r$ is
\[
\eta\int_0^r(A_t+\rho u)\,\dd u
=\eta\lrb{A_tr+\tfrac12\rho r^2}.
\]
Equating this expression to an $\operatorname{Exp}(1)$ variate gives the next
opening by a quadratic solve; closing waits are
$\operatorname{Exp}(\gamma)$. In the surging model, the opening wait is simply
$\operatorname{Exp}(\nu)$. The fixed-gate flows also solve exactly, so both
models admit event-driven simulation without time-step error in the switching.
Their analysis nevertheless begins with the joint generator
\cref{m:eq:coupledgenerator}: the countable-state methods of
\mextref{m:sec:ctmethods}{Section~1.3.2} do not by themselves close the
continuous-state transport or the mixed moment $\ex{G_tA_t}$.

\FloatBarrier
